% ============================================================================
% Dutch UI and structural strings for the One Chemistry Book series.
% Loaded by styles/onechemistry.sty when \booklang is nl.
% ============================================================================

% Theorem / statement environment titles
\newcommand{\omnameDefinition}{Definitie}
\newcommand{\omnameTheorem}{Stelling}
\newcommand{\omnameProposition}{Propositie}
\newcommand{\omnameLemma}{Lemma}
\newcommand{\omnameCorollary}{Gevolg}
\newcommand{\omnameMethod}{Methode}
\newcommand{\omnameExample}{Voorbeeld}
\newcommand{\omnameNotation}{Notatie}
\newcommand{\omnameRemark}{Opmerking}
\newcommand{\omnameExercise}{Oefening}
\newcommand{\omnameExercises}{Oefeningen}
\newcommand{\omnameProblem}{Probleem}
\newcommand{\omnameProblems}{Problemen}
\newcommand{\omnameProof}{Bewijs}

% Admitted results and solutions appendix headers
\newcommand{\omadmittedtext}{Op dit niveau zonder bewijs aangenomen.}
\newcommand{\omsolutionof}{Oplossing van}
\newcommand{\ompage}{pagina}
\newcommand{\omnameGotoSolution}{Oplossing p.}

% Misc text macros
\newcommand{\st}{\text{ zodat }}

% Cover / title page
\newcommand{\bookmaintitle}{One Chemistry Book}
\newcommand{\booktagline}{Het scheikundeboek om ze allen te regeren}
\newcommand{\bookdescription}{Scheikunde van het eerste schooljaar tot het einde van de bachelor}
\newcommand{\bookcontributors}{De bijdragers van het One Chemistry Book}
\newcommand{\bookversionlabel}{Versie}

% Colophon
\newcommand{\omcolophonsource}{Broncode, issues en bijdragen:}
\newcommand{\omcolophonlicense}{Tenzij anders vermeld, valt dit boek onder de licentie Creative Commons CC BY-NC-SA 4.0:}
\newcommand{\omcolophoncredit}{Bronvermelding: One Chemistry Book, One Course (one-course.com).}

% Chemistry boxes (styles/onechemistry.sty): the recall box that opens a
% chapter building on another one, the lab box, the history box, the safety
% box. Placeholder wording by the coordinator (2026-10-04); the Book 1
% `nl` agent owns it and settles the register (the recall title addresses
% the reader in the same register as the book's exercise stems).
\newcommand{\omnameRecall}{Wat je al weet}
\newcommand{\omnameInTheLab}{In het laboratorium}
\newcommand{\omnameHistory}{Geschiedenis}
\newcommand{\omnameSafety}{Veiligheid}

% Solutions appendix part title
\newcommand{\omsolutionspart}{Oplossingen van de oefeningen}

% Part titles (Jaren 1–12; parallel to English Grade N)
\@namedef{omstr@part.grade1}{Jaar 1 (6--7 jaar)}
\@namedef{omstr@part.grade2}{Jaar 2 (7--8 jaar)}
\@namedef{omstr@part.grade3}{Jaar 3 (8--9 jaar)}
\@namedef{omstr@part.grade4}{Jaar 4 (9--10 jaar)}
\@namedef{omstr@part.grade5}{Jaar 5 (10--11 jaar)}
\@namedef{omstr@part.grade6}{Jaar 6 (11--12 jaar)}
\@namedef{omstr@part.grade7}{Jaar 7 (12--13 jaar)}
\@namedef{omstr@part.grade8}{Jaar 8 (13--14 jaar)}
\@namedef{omstr@part.grade9}{Jaar 9 (14--15 jaar)}
\@namedef{omstr@part.grade10}{Jaar 10 (15--16 jaar)}
\@namedef{omstr@part.grade11}{Jaar 11 (16--17 jaar)}
\@namedef{omstr@part.grade12}{Jaar 12 (17--18 jaar)}

% Part titles (Bachelorjaren 1–3; parallel to English Bachelor Year N)
\@namedef{omstr@part.bachelor1}{Bachelorjaar 1 (18--19 jaar)}
\@namedef{omstr@part.bachelor2}{Bachelorjaar 2 (19--20 jaar)}
\@namedef{omstr@part.bachelor3}{Bachelorjaar 3 (20--21 jaar)}

% Plural names + \omnameChapter, used by the \crefname declarations in
% onechemistry.sty. Without these, cleveref falls back to its English built-ins.
\newcommand{\omnameDefinitions}{Definities}
\newcommand{\omnameTheorems}{Stellingen}
\newcommand{\omnamePropositions}{Proposities}
\newcommand{\omnameLemmas}{Lemma's}
\newcommand{\omnameCorollarys}{Gevolgen}
\newcommand{\omnameMethods}{Methoden}
\newcommand{\omnameExamples}{Voorbeelden}
\newcommand{\omnameNotations}{Notaties}
\newcommand{\omnameRemarks}{Opmerkingen}
\newcommand{\omnameChapter}{Hoofdstuk}
\newcommand{\omnameChapters}{Hoofdstukken}
% Section, for \cref{sec:…} (2026-10-04: it had no \crefname, so cleveref
% printed its English built-in "Section" in every edition).
\newcommand{\omnameSection}{Paragraaf}
\newcommand{\omnameSections}{Paragrafen}

% LaTeX's own structural names. babel would normally localise these, but it is
% optional here (see onechemistry.sty).
% siunitx and cleveref keep their English conjunctions unless told otherwise:
% \qtyrange printed "20 to 480 pF" and \cref{a,b} printed "Hoofdstukken 25 and
% 29" in the middle of Dutch prose. (The other language files have the same gap.)
% cleveref installs its conjunctions at \begin{document}, so redefine them there.
% siunitx typesets range phrases and list separators in MATH MODE, where a
% bare string loses its spaces (\qtyrange inside $...$ printed "1a100") and a
% non-ASCII character is worse: "a" expands to \`a, which LaTeX reports as
% "Command \` invalid in math mode" and drops from the page. siunitx's own
% defaults are wrapped in \text{} for exactly this reason; these overrides
% must be too. Found on Biology Book 3 fr, 2026-09-05, where the book's 56
% \qtyrange sites made it visible; it was latent in every language edition of
% every book in this repo.
\sisetup{range-phrase={\text{ tot }}, list-pair-separator={\text{ en }},
         list-final-separator={\text{ en }}}
\AtBeginDocument{%
  \renewcommand{\crefpairconjunction}{ en }%
  \renewcommand{\creflastconjunction}{ en }%
  \renewcommand{\crefrangeconjunction}{ tot }%
  % a \cref spanning two label types builds a group list, which has its
  % own pair of conjunctions: "Stellingen 2.14 en 2.18 and Lemma 2.17".
  \renewcommand{\crefpairgroupconjunction}{ en }%
  \renewcommand{\creflastgroupconjunction}{ en }%
}

% \today is English without babel (which is optional here), so the cover date
% read "July 26, 2026" on a Dutch book. Dutch writes it as "26 juli 2026",
% month lowercase.
\renewcommand{\today}{\number\day~\ifcase\month\or januari\or februari\or
  maart\or april\or mei\or juni\or juli\or augustus\or september\or oktober\or
  november\or december\fi\space\number\year}

\renewcommand{\chaptername}{Hoofdstuk}
\renewcommand{\partname}{Deel}
\renewcommand{\contentsname}{Inhoudsopgave}
\renewcommand{\appendixname}{Bijlage}
\renewcommand{\indexname}{Index}

% Periodic-table legend (\omperiodictable[legend], styles/onechemistry.sty).
% Coordinator wording, 2026-10-04 (it used to be hard-coded English in the
% style file, so no edition could translate it); the Book 1 agent of this
% language checks it against the book's own terminology.
\newcommand{\omnamePTmetal}{metalen}
\newcommand{\omnamePTmetalloid}{metalloïden}
\newcommand{\omnamePTnonmetal}{niet-metalen}
\newcommand{\omnamePTs}{s-blok}
\newcommand{\omnamePTp}{p-blok}
\newcommand{\omnamePTd}{d-blok}
\newcommand{\omnamePTf}{f-blok}
\newcommand{\omnamePTalkali}{alkalimetalen}
\newcommand{\omnamePTalkaline}{aardalkalimetalen}
\newcommand{\omnamePThalogen}{halogenen}
\newcommand{\omnamePTnoble}{edelgassen}
